% !TEX root = ../main.tex
\chapter{The Ross Recovery Theorem}\label{ch:ross}

\begin{sources}
\small
\begin{tabularx}{\linewidth}{@{}l l L@{}}
\toprule
\textbf{Lecture} & \textbf{Speaker} & \textbf{Content}\\
\midrule
2013 L25 & P.~Carr (Morgan Stanley), with J.~Yu & Guest lecture ``Can We Recover?'': Arrow--Debreu
security prices and market beliefs; Ross's Recovery Theorem for finite-state Markov chains via the
Perron--Frobenius theorem; a preference-free re-derivation using a change of numeraire (Long's numeraire
portfolio) for a time-homogeneous diffusion on a bounded state space; successes and failures for
unbounded state spaces (Black--Scholes fails, Cox--Ingersoll--Ross succeeds).\\
\bottomrule
\end{tabularx}

\smallskip
\textbf{Files.} \file{lecnote25} (50 slides, S.~Ross's collaborators P.~Carr and J.~Yu, Morgan Stanley).
There is no 2024 counterpart: the topic is not in the 2024 syllabus.
\textbf{Problem sets.} None (guest lecture, no associated problem-set questions in either year).
\textbf{Not covered / not published.} The slides' Part~I (conditioning and path-dependent
Arrow--Debreu prices, multi-period Markovian consistency) is background that the lecture itself
skipped over quickly (``I'm going to skip these slides''); it is summarised briefly below for
completeness but not derived in full. Ross's original paper (\emph{The Recovery Theorem}, published in
the \emph{Journal of Finance} in 2015, then only forthcoming) is not itself a course file; statements
attributed to it are as quoted or paraphrased on the slides and in the lecture.
\end{sources}

\section*{Motivation}
Every derivative price already encodes a probability measure: the \emph{risk-neutral} measure $\Q$ of
\cref{sec:oneperiod}, under which discounted prices are martingales. But $\Q$ is famously not the market's
actual view of the future --- it is contaminated by risk aversion, so that, say, the risk-neutral chance of
a market crash is typically far larger than anyone's honest assessment of its physical probability. For
decades the going wisdom was that untangling the two required an extra ingredient that no one could agree
on: a representative investor's risk-aversion parameter, supplied from outside the data. Ross's 2011
insight was that a single structural assumption on how risk aversion enters --- \emph{transition
independence} of the pricing kernel --- is enough to pin down the real-world probabilities uniquely, with
no free parameter to calibrate, provided the underlying is a finite-state, time-homogeneous Markov chain
and the matrix of state prices is positive. The mechanism is one this course has already built:
\cref{sec:pf}'s Perron--Frobenius theorem, applied to the matrix of Arrow--Debreu prices instead of a
Markov transition matrix. This chapter states the assumptions precisely, proves the recovery theorem in
full using \cref{thm:pf}, verifies the construction on an explicit numerical example, and closes Part~II of
these notes with a look at what happens (and what is still open) beyond the finite-state setting.

% ==================================================================
\section{Three probability measures: \texorpdfstring{$P$, $\Q$, $R$}{P, Q, R}}\label{ch25:sec-pqr}

Fix a single source of uncertainty $X$ (an index level, a short rate, an exchange rate) and consider
Arrow--Debreu (AD) securities written on it: a claim paying \$1 if $X$ lands in a given state and nothing
otherwise, i.e.\ exactly the state prices $\psi_i$ of \cref{def:onepmodel}, now indexed by both the
\emph{current} state and the terminal state \ts{13}{25}{12:49}. Three probability measures matter here
\ts{13}{25}{5:22}:
\begin{itemize}
  \item $P$, the \emph{physical} (objective, ``God's-eye'') probability measure: the true frequency with
        which $X$ visits each state. No one observes $P$ directly.
  \item $\Q$, the \emph{risk-neutral} measure of \cref{thm:ftap}: it exists under no arbitrage and
        satisfies $A_0=\beta\,\E_\Q[\bm C_T]$ for every replicable claim, but it is not a measure of frequency
        --- it is a repackaging of AD prices, normalised to integrate to one, and under it every asset grows
        at the risk-free rate \ts{13}{25}{9:37}.
  \item $R$, the \emph{recovered} (or \emph{natural}) measure: what one gets by trying to strip $\Q$ of its
        risk-aversion and time-value-of-money content and recover a measure of actual frequencies. Ross
        calls this the market's ``natural probabilities''. $R$ need not equal $P$ --- the market can be
        wrong (2005 mortgage-backed pricing being Carr's stock example) --- so $R$ is honestly only the
        market's \emph{belief} about $P$; a believer in market efficiency is free to set $R=P$
        \ts{13}{25}{10:40}.
\end{itemize}

\begin{intuition}[Why $\Q$ is not a frequency]
A risk-neutral probability of $40\%$ for ``the index is up'' means only that a security paying \$1 in that
event costs \$0.40 today, i.e.\ it is a \emph{price}, not a count of outcomes. Two effects separate this
price from the true chance of the event: discounting (a sure \$1 tomorrow is worth less than \$1 today,
even with no uncertainty at all) and risk aversion (a payoff arriving exactly when the rest of one's wealth
is doing badly is worth more per dollar than one arriving in good times, so its price is inflated relative
to its frequency). Recovering $R$ from $\Q$ means undoing both effects \ts{13}{25}{42:38}.
\end{intuition}

% ==================================================================
\section{The pricing matrix and the pricing kernel}\label{ch25:sec-kernel}

Let $X$ be a Markov chain on a finite state space $\{1,\dots,n\}$, and write $A(x,y)>0$ for the price, at
state $x$ today, of an AD security paying \$1 if $X$ is at state $y$ one period later \ts{13}{25}{40:28}.
Collecting these into a matrix $A\in\R^{n\times n}$, no-arbitrage consistency across horizons forces $X$'s
one-step transition structure under $A$ to be genuinely Markovian and, if desired, time-homogeneous, i.e.\
$A$ does not depend on calendar time (\emph{Appendix, Part~I of the slides}, sketched only: with three
dates $0,1,2$, path-dependent AD prices $A_{jk|i}$ factor as $A_{jk|i}=A_{ji}\cdot A_{k|ji}$, and Markov
means the one-step-ahead conditional price $A_{k|ji}=A_{k|j}$ does not depend on the earlier state $i$
\ts{13}{25}{39:23}).

\begin{coursenote}[Notation clash with \cref{ch:linalg}]
The letter $A$ here denotes the $n\times n$ matrix of \emph{one-step Arrow--Debreu prices between states of
a single Markov chain}, following Ross's and Carr's notation; it is unrelated to the $m\times n$ payoff
matrix $A$ of \cref{def:onepmodel} (assets in columns, states in rows). Both are conventional in their
respective sources and are kept as in the lecture material.
\end{coursenote}

Write $p(x,y)$ for the (unknown) \emph{natural} one-step transition probabilities under $R$: $p(x,y)\ge0$,
$\sum_yp(x,y)=1$ for every $x$. The \emph{pricing kernel} is the entrywise ratio
\begin{equation}\label{ch25:eq-kernel}
  \varphi(x,y) \;=\; \frac{A(x,y)}{p(x,y)},
\end{equation}
capturing exactly the two contaminating effects: time value of money and risk aversion
\ts{13}{25}{47:03}. If interest rates were zero and investors risk-neutral, $\varphi\equiv1$ and $A=p$;
in general $\varphi$ is an unknown positive function of two variables, and the $n^2$ equations \eqref{ch25:eq-kernel} involve
$2n^2-n$ unknowns (the $n^2$ values of $\varphi$ and the $n^2-n$ free entries of $p$, whose rows sum to one),
far too many to pin down $p$ --- this is why decontamination was thought hopeless before Ross
\ts{13}{25}{47:03}\ts{13}{25}{58:40}.

\begin{definition}[Transition independence, Ross's assumption]\label{ch25:def-ti}
The pricing kernel is \emph{transition-independent} if it factors as
\begin{equation}\label{ch25:eq-ti}
  \varphi(x,y) \;=\; \delta\,\frac{z(y)}{z(x)}, \qquad x,y=1,\dots,n,
\end{equation}
for some scalar $\delta>0$ and some function $z:\{1,\dots,n\}\to(0,\infty)$.
\end{definition}

In a world with a representative agent and additively time-separable utility, $z(x)=U'(c(x))$ (marginal
utility of consumption at state $x$) and $\delta$ is the subjective discount factor, so that
$\varphi(x,y)=\delta\,U'(c(y))/U'(c(x))$: risk aversion enters only through how marginal utility differs
\emph{between} the two states, not through the path or timing otherwise \ts{13}{25}{49:10}. Carr's
restatement, avoiding the representative-agent language, is purely dimension-counting
(\cref{ch25:rem-dof}); either reading is the same assumption \eqref{ch25:eq-ti}.

\begin{remark}[Degrees of freedom]\label{ch25:rem-dof}
Without \eqref{ch25:eq-ti}, $\varphi$ has $n^2$ free values. Assumption~\eqref{ch25:eq-ti} reduces this to
$n+1$: the $n$ values $z(1),\dots,z(n)$ (only their ratios matter, so really $n-1$ free numbers) plus
$\delta$. For $n=10$ this is a drop from $100$ to $11$ unknowns, which is what makes identification from
$A$ alone possible \ts{13}{25}{59:42}.
\end{remark}

% ==================================================================
\section{The Recovery Theorem}\label{ch25:sec-thm}

\begin{theorem}[Ross's Recovery Theorem]\label{ch25:thm-ross}
Let $A\in\R^{n\times n}$ be a matrix of one-step AD prices for a time-homogeneous Markov chain, with all
entries $A(x,y)>0$. Under transition independence \eqref{ch25:eq-ti}, there is a unique triple $(\delta,z,p)$
with $\delta>0$, $z\in\R^n$, $z>0$ (up to an overall positive scale factor, which cancels out of $p$ and
$\varphi$), and $p$ a row-stochastic matrix with $p(x,y)>0$, consistent with \eqref{ch25:eq-kernel}--%
\eqref{ch25:eq-ti} and the given $A$. Explicitly, $\delta$ is the Perron--Frobenius eigenvalue of $A$,
$w(x)=1/z(x)$ is (up to scale) its positive eigenvector, and
\begin{equation}\label{ch25:eq-precover}
  p(x,y) \;=\; \frac{1}{\delta}\,A(x,y)\,\frac{w(y)}{w(x)}, \qquad x,y=1,\dots,n.
\end{equation}
\end{theorem}

\begin{proof}
Write $D_z=\diag(z(1),\dots,z(n))$. Assumption \eqref{ch25:eq-ti} in \eqref{ch25:eq-kernel} reads
$A(x,y)=\delta\,p(x,y)\,z(y)/z(x)$, i.e.\ $z(x)A(x,y)=\delta\,p(x,y)z(y)$, or in matrix form
\begin{equation}\label{ch25:eq-matrixform}
  D_zA \;=\; \delta\,p\,D_z .
\end{equation}
Let $w=D_z^{-1}\bm1=(1/z(1),\dots,1/z(n))\T$, where $\bm1=(1,\dots,1)\T$. Multiply \eqref{ch25:eq-matrixform}
on the right by $w=D_z^{-1}\bm1$:
\[
  D_zAw \;=\; \delta\,p\,D_zD_z^{-1}\bm1 \;=\; \delta\,p\,\bm1 \;=\; \delta\bm1,
\]
using that $p$ is row-stochastic ($p\bm1=\bm1$). Multiplying through by $D_z^{-1}$,
\begin{equation}\label{ch25:eq-eigen}
  Aw \;=\; \delta\,D_z^{-1}\bm1 \;=\; \delta\,w .
\end{equation}
So $w$ is an eigenvector of $A$ with eigenvalue $\delta$, and $w>0$ because $z>0$.

\emph{Existence and uniqueness of $(\delta,w)$.} $A$ has all entries strictly positive, so
\cref{thm:pf} (Perron--Frobenius, proved in \cref{ch:linalg}) applies directly: $A$ has a real eigenvalue
$\lambda_0>0$ strictly larger in modulus than every other eigenvalue, with a positive eigenvector unique up
to positive scaling, and no other eigenvalue has a positive eigenvector (every other eigenvector has entries
of mixed sign or is complex). Equation~\eqref{ch25:eq-eigen} exhibits $w>0$ as a positive eigenvector of
$A$, so by \cref{thm:pf}(iii) it must be \emph{the} Perron eigenvector, unique up to scale, and $\delta$
must be \emph{the} Perron eigenvalue $\lambda_0$: no other pair $(\delta',w')$ with $w'>0$ solving
\eqref{ch25:eq-eigen} can exist. This proves $\delta$ and $z=1/w$ (up to the same overall scale) are unique.

\emph{Recovering $p$.} Solve \eqref{ch25:eq-matrixform} for $p$: $p=\delta^{-1}D_zAD_z^{-1}$, i.e.\ entrywise
$p(x,y)=\delta^{-1}z(x)A(x,y)/z(y)=\delta^{-1}A(x,y)w(y)/w(x)$, which is \eqref{ch25:eq-precover}. Positivity
$p(x,y)>0$ is immediate from $A,w,\delta>0$. Row-stochasticity is automatic and does not need to be imposed
separately: by \eqref{ch25:eq-eigen},
\[
\begin{gathered}
  \sum_{y=1}^np(x,y) \;=\; \frac{1}{\delta\,w(x)}\sum_{y=1}^nA(x,y)w(y) \;=\; \frac{(Aw)(x)}{\delta\,w(x)} \\
  \;=\; \frac{\delta\,w(x)}{\delta\,w(x)} \;=\; 1,
\end{gathered}
\]
for every $x$. Since $(\delta,w)$ are the unique Perron data of $A$, $p$ in \eqref{ch25:eq-precover} is the
unique matrix consistent with \eqref{ch25:eq-ti} and the given $A$.
\end{proof}

\begin{coursenote}[What the lecture actually showed]
Carr states Ross's Theorem~1 essentially verbatim (``if the pricing matrix $A$ is positive or irreducible,
then there exists a unique solution... via the Perron--Frobenius theorem'') but says explicitly that he has
no slide on \emph{how} the $n+1$ unknowns are actually computed, referring the audience to Ross's paper
\ts{13}{25}{1:01:49}. The derivation above (via \eqref{ch25:eq-matrixform}--\eqref{ch25:eq-eigen}) is a
reconstruction, not shown on the slides or in the transcript, built from the definitions Carr does give
(the pricing kernel \eqref{ch25:eq-kernel}, transition independence \eqref{ch25:eq-ti}, and the explicit
statement that Perron--Frobenius is the tool used). The theorem also holds, with $|\lambda|\le\delta$ rather
than strict inequality, for \emph{irreducible non-negative} $A$ (some $A(x,y)$ may vanish): this is the
version of \cref{thm:pf} noted in \cref{ch:linalg} as the one needed for Markov chains, and it is the
version Ross actually states.
\end{coursenote}

\begin{finance}[Reading the recovered kernel]
Once $p$ is known, \eqref{ch25:eq-ti} gives back the pricing kernel itself,
$\varphi(x,y)=\delta\,z(y)/z(x)=\delta\,w(x)/w(y)$, fully decontaminated: $\delta$ is the (state-independent)
discount factor and $w(x)/w(y)$ isolates the risk-aversion premium of moving from $x$ to $y$. If the recovered marginal utility
$z=1/w$ is decreasing in the state index (i.e.\ $w$ is increasing) when the index orders states from
``bad'' to ``good'', a transition into a worse state costs relatively more than a fair bet on its recovered
probability alone would suggest --- exactly the insurance-value story of \cref{ch25:sec-pqr}. Unlike the
one-period model of \cref{ch:linalg}, where \emph{any} positive $\bm\psi$ with $A_0=\bm\psi\T A$ is an
admissible state-price vector when the market is incomplete, here positivity of the whole matrix $A$ plus
\eqref{ch25:eq-ti} is what removes all the remaining freedom in one step, without ever specifying a utility
function numerically.
\end{finance}

% ==================================================================
\section{Worked example and numerical verification}\label{ch25:sec-example}

\begin{example}[Recovering a 3-state chain]\label{ch25:ex-3state}
Take three states (say, three ranges of a 1-year index return: down, flat, up) with the ``true'' natural
transition matrix
\[
  p = \begin{pmatrix} 0.6 & 0.3 & 0.1 \\ 0.2 & 0.5 & 0.3 \\ 0.1 & 0.4 & 0.5 \end{pmatrix},
\]
a discount factor $\delta=0.97$, and marginal utilities $z=(1.0,\,1.6,\,2.5)$ (higher in the ``up'' state,
i.e.\ increasing marginal utility of consumption in the better state --- an unusual but not excluded case,
matching Ross's remark that the recovered kernel could be increasing in $y$ \ts{13}{25}{55:28}). By
\eqref{ch25:eq-ti}--\eqref{ch25:eq-kernel}, the (only observable, in this thought experiment) AD price
matrix is
\[
  A(x,y) = \delta\, p(x,y)\, \frac{z(y)}{z(x)}
  = \begin{pmatrix} 0.5820 & 0.4656 & 0.2425 \\ 0.1213 & 0.4850 & 0.4547 \\ 0.0388 & 0.2483 & 0.4850 \end{pmatrix}.
\]
Given only $A$, \cref{ch25:thm-ross} recovers $\delta$ as its Perron eigenvalue and $w=1/z$ as its Perron
eigenvector (fixing the scale by, say, $w(1)=1$); \eqref{ch25:eq-precover} then reconstructs $p$ exactly,
with no knowledge of $z$, $\delta$, or $p$ used in the recovery step itself.
\end{example}

\begin{casestudy}[Numerical check, Python]
The following reconstructs \cref{ch25:ex-3state}: build $A$ from $(p,\delta,z)$, forget $(p,\delta,z)$, and
recover them from $A$ alone via \texttt{numpy.linalg.eig}.
\begin{lstlisting}[language=Python]
import numpy as np

P = np.array([[0.6, 0.3, 0.1],
              [0.2, 0.5, 0.3],
              [0.1, 0.4, 0.5]])
delta, z = 0.97, np.array([1.0, 1.6, 2.5])
A = delta * P * (z[None, :] / z[:, None])     # build the AD price matrix

eigvals, eigvecs = np.linalg.eig(A)           # forget P, delta, z; recover them
i = np.argmax(eigvals.real)                   # Perron eigenvalue = largest one
lam0, w = eigvals[i].real, np.abs(eigvecs[:, i].real)
w /= w[0]                                     # fix the free positive scale

P_recovered = (1 / lam0) * A * (w[None, :] / w[:, None])
print(lam0, w)                 # -> 0.97, [1.0, 0.625, 0.4] = 1/z (up to scale)
print(np.max(np.abs(P_recovered - P)))         # -> ~4e-16
\end{lstlisting}
The recovered eigenvalue is $\lambda_0=0.97=\delta$ to machine precision, the eigenvector is
$w=(1,\,0.625,\,0.4)=1/z$ up to the chosen scale, and $P_{\mathrm{recovered}}$ matches the true $p$ to
within $4\times10^{-16}$ (floating-point roundoff): \cref{ch25:eq-precover} reproduces the natural
transition matrix exactly from the price matrix alone.
\end{casestudy}

\begin{intuition}[Why positivity of $A$ is doing all the work]
Nothing in the recovery step uses probabilistic reasoning about $p$ directly; it is pure linear algebra on
$A$. The only place probability enters is in guaranteeing, via no-arbitrage, that $A$ is a matrix of
strictly positive prices in the first place, and in the modelling assumption \eqref{ch25:eq-ti} that
restricts what $A$ can look like. Once both hold, the Perron eigenvector is forced to be positive and simple
by \cref{thm:pf}, and positivity of the reconstructed $p$ and its automatic row-sum-one property follow for
free --- there is no need to separately check that the answer is a sensible probability matrix.
\end{intuition}

% ==================================================================
\section{Beyond finite-state chains: change of numeraire}\label{ch25:sec-diffusion}

Ross's own proof works in the finite-state, discrete-time setting above. Carr and Yu's contribution is a
preference-free re-derivation, avoiding the representative-agent language, that extends the same conclusion
to a continuous-state diffusion. The key substitute for the discrete pricing kernel is a change of
numeraire.

\begin{definition}[Numeraire, numeraire portfolio]\label{ch25:def-numeraire}
A \emph{numeraire} is a self-financing portfolio whose value is strictly positive at all times (a swap can
be zero or negative, hence disqualified; a stock or a money market account qualifies, as long as it never
hits zero) \ts{13}{25}{25:20}. Long (1990) showed that, given no arbitrage among a money market account and
$n$ risky assets, there is always a self-financing portfolio of them, of value $L_t>0$, the \emph{numeraire
portfolio}, such that every relative price $S^i_t/L_t$ is an $R$-martingale --- i.e.\ expressing P\&L in
units of $L$ turns the (unobservable) natural measure $R$ into the measure under which discounted prices
behave like ordinary risk-neutral martingales \ts{13}{25}{1:16:14}. The same object is also called the
\emph{growth-optimal portfolio} (Kelly, 1956: over an infinite horizon it has the largest mean growth rate
among all portfolios, at the price of being quite risky along the way) or the \emph{natural numeraire}
\ts{13}{25}{19:04}.
\end{definition}

Assume $X$ is a univariate, time-homogeneous diffusion under $\Q$,
\[
  dX_t = b^\Q(X_t)\,dt + a(X_t)\,dW_t,
\]
on a \emph{bounded} state space $(\ell,u)$, driving the short rate $r_t=r(X_t)$ and the value of $n$ risky
assets, all with known $\Q$-coefficients; and that the numeraire portfolio's value is a function
$L_t=L(X_t,t)$ solving $dL_t/L_t=r(X_t)\,dt+\sigma_L(X_t)\,dW_t$ for an unknown volatility function
$\sigma_L$ \ts{13}{25}{1:17:17}. Self-financing forces $L$ to solve the pricing PDE
$\partial_tL+\tfrac12a^2\partial_x^2L+b^\Q\partial_xL=rL$; writing $\sigma_L(x)=a(x)\,\partial_x\ln L(x,t)$
(from It\^o's formula) and separating variables $L(x,t)=\pi(x)e^{\lambda t}$ turns this PDE into a
\emph{regular Sturm--Liouville problem}
\begin{equation}\label{ch25:eq-sl}
  \frac{a^2(x)}{2}\pi''(x) + b^\Q(x)\pi'(x) - r(x)\pi(x) \;=\; -\lambda\,\pi(x), \qquad x\in(\ell,u),
\end{equation}
with separated boundary conditions at $\ell,u$ \ts{13}{25}{1:25:39}.

\begin{finance}[Sturm--Liouville as the continuum Perron--Frobenius]
Equation~\eqref{ch25:eq-sl} is the exact continuous-state analogue of the eigenvalue problem
\eqref{ch25:eq-eigen}: Sturm--Liouville theory guarantees a smallest eigenvalue $\lambda_0$, strictly below
every other eigenvalue, with an associated eigenfunction $\pi_0(x)>0$ everywhere and unique up to positive
scaling, while every other eigenfunction changes sign at least once \ts{13}{25}{1:26:40} --- precisely the
conclusions (i)--(iii) of \cref{thm:pf} carried over from finitely many states to a continuum of them. Once
$\pi_0$ and $\lambda_0$ are known (numerically, in general), $L(x,t)=\pi_0(x)e^{\lambda_0t}$ is known up to
scale, hence so is $\sigma_L(x)=a(x)\,\partial_x\ln\pi_0(x)$, and Girsanov's theorem then gives the $R$-drift
of $X$ as $b^\Q(x)+a(x)\sigma_L(x)=b^\Q(x)+a^2(x)\,\partial_x\ln\pi_0(x)$ and the $R$-dynamics of every asset in the set \ts{13}{25}{1:19:23}.
This is the ``$11$ unknowns'' of \cref{ch25:rem-dof} replaced by one unknown function $\pi_0$ and one
unknown scalar $\lambda_0$.
\end{finance}

\begin{coursenote}[Boundedness is essential, and still an open problem beyond it]
The construction needs $X$'s state space to be an interval $(\ell,u)$ with $\ell,u$ finite (Carr is candid
that some economists object to bounding, say, the S\&P~500 above, and that he personally has no problem
with a bound of \$20 trillion \ts{13}{25}{27:29}). For an \emph{unbounded} state space the argument needs a
Hilbert-space (square-integrability) version of Sturm--Liouville theory, and it does not always go through:
it \emph{fails} for the standard geometric Brownian motion of Black--Scholes (so the natural measure of a
stock price cannot be recovered this way from stock option prices under these assumptions), but it
\emph{succeeds} for the Cox--Ingersoll--Ross short-rate model, recovering an R-dynamics that is again CIR
but with a higher mean-reversion speed. More generally it succeeds for affine diffusions (Linetsky and a
student, cited but not derived in the lecture). At the time of the 2013 lecture, necessary and sufficient
conditions for recovery on an unbounded state space were an open research question
\ts{13}{25}{1:27:41}.
\end{coursenote}

% ==================================================================
\begin{keyformulas}
\begin{align*}
&\text{Pricing kernel:} && \varphi(x,y)=A(x,y)/p(x,y)\\
&\text{Transition independence:} && \varphi(x,y)=\delta\,z(y)/z(x)\iff D_zA=\delta\,p\,D_z\\
&\text{Recovery ($w=1/z$):} && Aw=\delta w \ \text{(Perron eigenpair of }A\text{, \cref{thm:pf})}\\
&\text{Recovered transitions:} && p(x,y)=\delta^{-1}A(x,y)\,w(y)/w(x),\quad\textstyle\sum_yp(x,y)=1\\
&\text{Diffusion analogue:} && \tfrac12a^2\pi''+b^\Q\pi'-r\pi=-\lambda\pi\ \text{(Sturm--Liouville)}\\
  &&& \sigma_L(x)=a(x)\partial_x\ln\pi_0(x),\quad b^R(x)=b^\Q(x)+a(x)\sigma_L(x)
\end{align*}
\end{keyformulas}

% ==================================================================
\section*{Exercises}
\addcontentsline{toc}{section}{Exercises}

\begin{exercise}[Recovering a 2-state chain by hand (not from the course)]\label{ch25:ex-x-2state}
Let $A=\bigl(\begin{smallmatrix}0.45&0.30\\0.20&0.40\end{smallmatrix}\bigr)$.
\begin{enumerate}[(a)]
  \item Find the Perron eigenvalue
      $\delta$ and a positive eigenvector $w$ of $A$ by solving the $2\times2$ characteristic equation directly
      (no software).
  \item Use \eqref{ch25:eq-precover} to recover $p$, and check it is row-stochastic.
  \item Recover
      $z=1/w$ up to scale and the pricing kernel $\varphi(x,y)=\delta z(y)/z(x)$; is $\varphi$ larger for a
      transition into state 1 or state 2? Which state carries the ``insurance'' premium?
\end{enumerate}
\end{exercise}

\begin{exercise}[Recovery is not just normalising rows (not from the course)]\label{ch25:ex-x-notnorm}
For the $A$ of \cref{ch25:ex-x-2state}, compute the naive row-normalisation
$\hat p(x,y)=A(x,y)/\sum_zA(x,z)$.
\begin{enumerate}[(a)]
  \item Verify $\hat p\ne p$ from \cref{ch25:ex-x-2state}(b).
  \item Explain, in
      terms of \eqref{ch25:eq-matrixform}, why row-normalising $A$ recovers the natural probabilities only in the
      special case $z\equiv\mathrm{const}$ (no risk aversion), and identify what $\hat p$ actually equals in
      general (a $\Q$-like object, not $R$).
\end{enumerate}
\end{exercise}

\begin{exercise}[Sensitivity of the recovered kernel (not from the course)]\label{ch25:ex-x-sensitivity}
Using the Python code of \cref{ch25:ex-3state}, perturb one entry of $A$ (e.g.\ multiply $A(1,3)$ by $1.05$)
while leaving the others fixed, keeping $A$ positive.
\begin{enumerate}[(a)]
  \item Recompute $\delta$, $w$ and $p$; by how much does
      each change relative to the unperturbed values?
  \item Which entries of $p$ are most sensitive to a change in
      $A(1,3)$, and does this match the intuition that Perron--Frobenius data depend on the \emph{whole} matrix,
      not just one entry?
  \item Repeat with a $5\times5$ random positive matrix and comment on whether the recovered
      $p$ remains close to row-stochastic under floating-point rounding alone.
\end{enumerate}
\end{exercise}

\begin{exercise}[Irreducible but not positive (not from the course)]\label{ch25:ex-x-irreducible}
Let $A=\bigl(\begin{smallmatrix}0&0.5&0.3\\0.4&0&0.6\\0.2&0.4&0\end{smallmatrix}\bigr)$ (zero diagonal, but
irreducible: every state reachable from every other in finitely many steps).
\begin{enumerate}[(a)]
  \item Check numerically that $A$
      still has a positive Perron eigenvector, consistent with the irreducible case of \cref{thm:pf} mentioned in
      \cref{ch:linalg}.
  \item Recover $p$ via \eqref{ch25:eq-precover} and check it is a valid (row-stochastic,
      non-negative) transition matrix with some zero entries.
  \item Explain financially what $A(x,x)=0$ means (no
      one-period AD security paying off in the same state you start in) and why this does not break recovery.
\end{enumerate}
\end{exercise}

\begin{exercise}[From the Sturm--Liouville problem to the discrete one (not from the course)]\label{ch25:ex-x-discretize}
Discretise \eqref{ch25:eq-sl} on a grid $x_1<\dots<x_n$ in $(\ell,u)$ using centred differences for
$\pi''(x_i)\approx(\pi_{i+1}-2\pi_i+\pi_{i-1})/h^2$ and $\pi'(x_i)\approx(\pi_{i+1}-\pi_{i-1})/(2h)$, with
absorbing (Dirichlet) boundary conditions $\pi_0=\pi_{n+1}=0$.
\begin{enumerate}[(a)]
  \item Write the resulting eigenvalue problem as
      $M\bm\pi=-\lambda\bm\pi$ for an explicit tridiagonal $M$ in terms of $a(x_i)$, $b^\Q(x_i)$, $r(x_i)$, $h$.
  \item For constant $a,b^\Q,r$, compute $M$'s eigenvalues in closed form and identify the smallest one and its
      eigenvector's sign pattern; check it matches \cref{thm:pf}'s conclusion that only the extremal eigenvalue
      has an everywhere-positive eigenvector.
  \item Implement this in Python for $a=0.2$, $b^\Q(x)=0.1(0.05-x)$,
      $r(x)=0.03$, on $(\ell,u)=(-0.1,0.2)$ with $n=50$ grid points, and plot $\pi_0$; is it unimodal?
\end{enumerate}
\end{exercise}

\begin{remark}[End of Part~II]
This closes Part~II of these notes. The recovery theorem was a fitting last stop: it leans on nearly every
tool built up along the way --- state prices and no-arbitrage (\cref{ch:linalg}), the risk-neutral measure
(\cref{ch:bs}), Markov chains and Perron--Frobenius (\cref{ch:linalg} again), and, in its diffusion form,
stochastic calculus and PDEs (\cref{ch:stochcalc,ch:sde}). The appendices that follow are practical rather
than mathematical: how to run the course's \textsf{R} material, and the logistics of both years of the
course.
\end{remark}
